\section{Discussion and Limitations}

\subsection{What the retained route establishes}

Stabilized NV/NV remains the higher-fidelity full-FP4 control. Its row
reference, E4M3 scale construction, and inverse correction are expensive on
the online path; removing them makes NV/NV distribution-dependent. NV/MX
chooses a different compromise. Its E8M0 scale represents small N32
probability blocks without a tensor-wide P factor and folds into the direct
log-domain code map, at the cost of coarser scale placement.

The resulting two-times-class advantage over BF16 is a kernel result, not the
four-times peak matrix ratio. BF16 FA4 already overlaps much of its non-matrix
work. FP4 shortens QK and PV but still loads FP32 scores, waits for K64-ready
pairs, publishes scales, updates online state, and stores output. The fixed-P
ceiling shows that the remaining P arithmetic alone is no longer enough to
close that gap.

The attribution is equally important. The two-query CTA, full 512-column TMEM
layout, warp specialization, deep K/V ring, and split P publication come from
HAO. This work contributes the mixed NV/MX representation, direct E2M1 code
map, SFU/FMA balance, represented denominator, sampled range guard, and its
underflow-safe evaluation.

\subsection{Numerical evidence and its limits}

Operator cosine and relative-$L_2$ are worse than HAO NV/FP8, while the tested
ViT, BERT, and ViT-MAE tasks show much smaller final changes. Residual paths,
normalization, later layers, and decision margins explain why these measures
need not move together. They do not prove that the approximation is safe for
every model or for training.

Wan provides a harder depth test. The corrected fast route is finite through
20 steps on calibration and held-out prompts, but its error grows more than
HAO NV/NV and NV/FP8. Layer-wise coefficient fitting also fails to transfer
reliably because a single-layer BF16-teacher proxy does not model interacting
errors in an all-FP4 trajectory. Future calibration should optimize a small
shared codebook under the end-task loss, ideally through QAT, and validate on
fresh prompts and processes.

The present suite is forward-only and noncausal. It excludes dynamic Q/K/V
quantization, folded-scale preparation, and K/V permutation cost. It does not
establish backward stability, causal decoding, GQA, variable-length support,
training convergence, language-model perplexity, long-context retrieval, or
full video-quality metrics. SageAttention3 and Attn-QAT provide the relevant
larger evaluation templates \cite{zhang2025sageattention3,zhang2026attnqat}.

\subsection{Hardware dependence}

This design is specific to SM100/SM103's TMEM and scaled-FP4 instruction
contracts. B300 improves favorable saturated shapes, but it has the same
256-KB TMEM capacity and the tested SKU exposes fewer SMs at a slightly lower
clock. One all-TMEM CTA still resides per SM, and broad use of its faster SFU
does not remove the P-to-PV dependency. A GPU with another score bank, K32
scaled PV, or direct compact-scale operands could prefer a deeper pipeline.
Conversely, SageAttention3's SM120 result informs the numerics but not this
TMEM schedule.
